% Part: proof-theory
% Chapter: normalization
% Section: permutations

\documentclass[../../../include/open-logic-section]{subfiles}

\begin{document}

\olfileid[id]{pt}{nor}{per}

\olsection{Konversi Permutasi}

Potongan dalam \Log{N1i} adalah segmen yang berakhir pada premis utama
!!a{elimination}. Salah satu kasus dalam bukti normalisasi ialah mengurangi
panjang potongan maksimal, dan dengan demikian panjang potongan
!!{proof}. Suatu potongan dengan panjang $>1$ dapat berakhir dalam banyak
cara, bergantung pada apakah kemunculan !!{formula}~$!A$ terakhir dalam
potongan itu merupakan kesimpulan dari \Elim{\lor} atau \Elim{\lexists},
dan bergantung pada inferensi !!{elimination} yang premis utamanya adalah
kemunculan tersebut.

Sebagai contoh, jika inferensi yang terlibat ialah \Elim{\lor} dan
\Elim{\land}, $!A \ident !B \lor !C$, dan terdapat
sub-!!{proof}~$\delta_1$ yang berakhir dengan:
\[
\AxiomC{}
\RightLabel{$\delta_2$}
\DeduceC{$!D \lor !E$}
\AxiomC{$\Discharge{!D}{x}$}
\RightLabel{$\delta_3$}
\DeduceC{$!B \land !C$}
\AxiomC{$\Discharge{!E}{y}$}
\RightLabel{$\delta_4$}
\DeduceC{$!B \land !C$}
\DischargeRule{\Elim{\lor}}{x\,y}
\TrinaryInfC{$!B \land !C$}
\RightLabel{\Elim{\land}[1]}
\UnaryInfC{$!B$}
\DisplayProof
\]
Kemunculan !!{formula} terakhir~$!A_n$ dalam potongan adalah kesimpulan
\Elim{\lor}, yakni $!B \land !C$ yang berada di bawah. Kemunculan
!!{formula} kedua dari belakang~$!A_{n-1}$ adalah salah satu premis minor
\Elim{\lor}.

Untuk mengurangi panjang potongan ini, kita dapat menukar urutan
(``mempermutasikan'') inferensi \Elim{\lor} dan \Elim{\land}. Artinya,
kita mengganti subbukti~$\delta_1$ dalam~$\delta$ dengan yang berikut:
\[
\AxiomC{}
\RightLabel{$\delta_2$}
\DeduceC{$!D \lor !E$}
\AxiomC{$\Discharge{!D}{x}$}
\RightLabel{$\delta_3$}
\DeduceC{$!B \land !C$}
\RightLabel{\Elim{\land}[1]}
\UnaryInfC{$!B$}
\AxiomC{$\Discharge{!E}{y}$}
\RightLabel{$\delta_4$}
\DeduceC{$!B \land !C$}
\RightLabel{\Elim{\land}[1]}
\UnaryInfC{$!B$}
\DischargeRule{\Elim{\lor}}{x\,y}
\TrinaryInfC{$!B$}
\DisplayProof
\]
Potongan-potongan yang sebelumnya berakhir pada premis \Elim{\land} di
bawah \Elim{\lor} kini berakhir pada premis salah satu inferensi
\Elim{\land} di atas premis-premis minor \Elim{\lor}; dengan kata lain,
panjang masing-masing berkurang~$1$.

Selain itu, setiap potongan yang seluruhnya terletak dalam $\delta_2$,
$\delta_3$, atau $\delta_4$, maupun seluruhnya di luar~$\delta_1$ dalam
$\delta$, tidak berubah: potongan tersebut masih melibatkan !!{formula}s
yang sama dan melewati inferensi yang sama, sehingga khususnya mempunyai
peringkat dan panjang yang sama dengan potongan bersesuaian dalam~$\delta$.
Jadi, peringkat potongan !!{proof} baru tidak meningkat, dan panjang
potongan !!{proof} baru lebih kecil daripada panjang potongan !!{proof} lama.

\begin{prob}
Setelah suatu konversi permutasi, panjang sekurang-kurangnya satu potongan
berkurang~$1$. Jika kita mempermutasikan \Elim{\land} melintasi
\Elim{\lor}, sebenarnya terdapat dua segmen potongan yang panjangnya
berkurang, sehingga dalam kasus ini panjang potongan !!{proof} berkurang
sekurang-kurangnya~$2$. Dapatkah panjang potongan !!a{proof} berkurang lebih
dari~$2$ melalui satu permutasi seperti itu? Dapatkah \emph{peringkat}
potongannya berkurang?
\end{prob}

Proses memindahkan !!a{elimination} ke atas aturan \Elim{\lor} atau
\Elim{\lexists} disebut \emph{konversi permutasi}. Pola bagi beberapa
pasangan aturan diberikan dalam \olref{tab:permutation-conversions}.
(Pola lainnya dibiarkan sebagai latihan.)

% The earlier table source is intentionally retained as a non-rendered source note.
% begin{table}
% begin{multline*}
% shoveleft{AxiomC{$!D lor !E$}
% AxiomC{$Discharge{!D}{x}$}
% shortDeduce
% DeduceC{$!B land !C$}
% AxiomC{$Discharge{!E}{y}$}
% shortDeduce
% DeduceC{$!B land !C$}
% DischargeRule{Elim{lor}}{x,y}
% TrinaryInfC{$!B land !C$}
% RightLabel{Elim{land}[1]}
% UnaryInfC{$!B$}
% DisplayProof}\
% shoveright{longrightarrow\ 
% AxiomC{$!D lor !E$}
% AxiomC{$Discharge{!D}{x}$}
% shortDeduce
% DeduceC{$!B land !C$}
% RightLabel{Elim{land}[1]}
% UnaryInfC{$!B$}
% AxiomC{$Discharge{!E}{y}$}
% shortDeduce
% DeduceC{$!B land !C$}
% RightLabel{Elim{land}[1]}
% UnaryInfC{$!B$}
% DischargeRule{Elim{lor}}{x,y}
% TrinaryInfC{$!B$}
% DisplayProof}
% \
% shoveleft{AxiomC{$!D lor !E$}
% AxiomC{$Discharge{!D}{x}$}
% shortDeduce
% DeduceC{$!B lif !C$}
% AxiomC{$Discharge{!E}{y}$}
% shortDeduce
% DeduceC{$!B lif !C$}
% DischargeRule{Elim{lor}}{x,y}
% TrinaryInfC{$!B lif !C$}
% AxiomC{$!B$}
% RightLabel{Elim{lif}}
% BinaryInfC{$!C$}
% DisplayProof}\
% shoveright{longrightarrow\ 
% AxiomC{$!D lor !E$}
% AxiomC{$Discharge{!D}{x}$}
% shortDeduce
% DeduceC{$!B lif !C$}
% AxiomC{$!B$}
% RightLabel{Elim{lif}}
% BinaryInfC{$!C$}
% AxiomC{$Discharge{!E}{y}$}
% shortDeduce
% DeduceC{$!B lif !C$}
% AxiomC{$!B$}
% RightLabel{Elim{lif}}
% BinaryInfC{$!C$}
% DischargeRule{Elim{lor}}{x,y}
% TrinaryInfC{$!C$}
% DisplayProof}
% \
% shoveleft{AxiomC{$!D lor !E$}
% AxiomC{$Discharge{!D}{x}$}
% shortDeduce
% DeduceC{$!B lif !C$}
% AxiomC{$Discharge{!E}{y}$}
% shortDeduce
% DeduceC{$!B lif !C$}
% DischargeRule{Elim{lor}}{x,y}
% TrinaryInfC{$!B lif !C$}
% AxiomC{$!B$}
% RightLabel{Elim{lif}}
% BinaryInfC{$!C$}
% DisplayProof}\
% shoveright{longrightarrow\ 
% AxiomC{$!D lor !E$}
% AxiomC{$Discharge{!D}{x}$}
% shortDeduce
% DeduceC{$!B lif !C$}
% AxiomC{$!B$}
% RightLabel{Elim{lif}}
% BinaryInfC{$!C$}
% AxiomC{$Discharge{!E}{y}$}
% shortDeduce
% DeduceC{$!B lif !C$}
% AxiomC{$!B$}
% RightLabel{Elim{lif}}
% BinaryInfC{$!C$}
% DischargeRule{Elim{lor}}{x,y}
% TrinaryInfC{$!C$}
% DisplayProof}
% \
% shoveleft{AxiomC{$!D lor !E$}
% AxiomC{$Discharge{!D}{x}$}
% shortDeduce
% DeduceC{$lforall[x][!B(x)]$}
% AxiomC{$Discharge{!E}{y}$}
% shortDeduce
% DeduceC{$lforall[x][!B(x)]$}
% DischargeRule{Elim{lor}}{x,y}
% TrinaryInfC{$lforall[x][!B(x)]$}
% RightLabel{Elim{lforall}}
% UnaryInfC{$!B(t)$}
% DisplayProof}\
% shoveright{longrightarrow\ 
% AxiomC{$!D lor !E$}
% AxiomC{$Discharge{!D}{x}$}
% shortDeduce
% DeduceC{$lforall[x][!B(x)]$}
% RightLabel{Elim{lforall}}
% UnaryInfC{$!B(t)$}
% AxiomC{$Discharge{!E}{y}$}
% shortDeduce
% DeduceC{$lforall[x][!B(x)]$}
% RightLabel{Elim{lforall}}
% UnaryInfC{$!B(t)$}
% DischargeRule{Elim{lor}}{x,y}
% TrinaryInfC{$!B(t)$}
% DisplayProof}
% end{multline*}
% caption{Some permutation conversions for Log{N1i}.}
% ollabel{tab:permutation-conversions}
% end{table}

{\small
\begin{longtable}{@{}l@{}}
\AxiomC{$!D \lor !E$}
\AxiomC{$\Discharge{!D}{x}$}
\shortDeduce
\DeduceC{$!B \land !C$}
\AxiomC{$\Discharge{!E}{y}$}
\shortDeduce
\DeduceC{$!B \land !C$}
\DischargeRule{\Elim{\lor}}{x\,y}
\TrinaryInfC{$!B \land !C$}
\RightLabel{\Elim{\land}[1]}
\UnaryInfC{$!B$}
\DisplayProof
$\longrightarrow$\\[6ex]
\quad\AxiomC{$!D \lor !E$}
\AxiomC{$\Discharge{!D}{x}$}
\shortDeduce
\DeduceC{$!B \land !C$}
\RightLabel{\Elim{\land}[1]}
\UnaryInfC{$!B$}
\AxiomC{$\Discharge{!E}{y}$}
\shortDeduce
\DeduceC{$!B \land !C$}
\RightLabel{\Elim{\land}[1]}
\UnaryInfC{$!B$}
\DischargeRule{\Elim{\lor}}{x\,y}
\TrinaryInfC{$!B$}
\DisplayProof
\\[10ex]
\AxiomC{$!D \lor !E$}
\AxiomC{$\Discharge{!D}{x}$}
\shortDeduce
\DeduceC{$!B \lif !C$}
\AxiomC{$\Discharge{!E}{y}$}
\shortDeduce
\DeduceC{$!B \lif !C$}
\DischargeRule{\Elim{\lor}}{x\,y}
\TrinaryInfC{$!B \lif !C$}
\AxiomC{$!B$}
\RightLabel{\Elim{\lif}}
\BinaryInfC{$!C$}
\DisplayProof
$\longrightarrow$\\[6ex]
\AxiomC{$!D \lor !E$}
\AxiomC{$\Discharge{!D}{x}$}
\shortDeduce
\DeduceC{$!B \lif !C$}
\AxiomC{$!B$}
\RightLabel{\Elim{\lif}}
\BinaryInfC{$!C$}
\AxiomC{$\Discharge{!E}{y}$}
\shortDeduce
\DeduceC{$!B \lif !C$}
\AxiomC{$!B$}
\RightLabel{\Elim{\lif}}
\BinaryInfC{$!C$}
\DischargeRule{\Elim{\lor}}{x\,y}
\TrinaryInfC{$!C$}
\DisplayProof
\\[10ex]
\AxiomC{$!D \lor !E$}
\AxiomC{$\Discharge{!D}{x}$}
\shortDeduce
\DeduceC{$!B \lif !C$}
\AxiomC{$\Discharge{!E}{y}$}
\shortDeduce
\DeduceC{$!B \lif !C$}
\DischargeRule{\Elim{\lor}}{x\,y}
\TrinaryInfC{$!B \lif !C$}
\AxiomC{$!B$}
\RightLabel{\Elim{\lif}}
\BinaryInfC{$!C$}
\DisplayProof
$\longrightarrow$\\[6ex]
\AxiomC{$!D \lor !E$}
\AxiomC{$\Discharge{!D}{x}$}
\shortDeduce
\DeduceC{$!B \lif !C$}
\AxiomC{$!B$}
\RightLabel{\Elim{\lif}}
\BinaryInfC{$!C$}
\AxiomC{$\Discharge{!E}{y}$}
\shortDeduce
\DeduceC{$!B \lif !C$}
\AxiomC{$!B$}
\RightLabel{\Elim{\lif}}
\BinaryInfC{$!C$}
\DischargeRule{\Elim{\lor}}{x\,y}
\TrinaryInfC{$!C$}
\DisplayProof
\\[10ex]
\AxiomC{$!D \lor !E$}
\AxiomC{$\Discharge{!D}{x}$}
\shortDeduce
\DeduceC{$\lforall[x][!B(x)]$}
\AxiomC{$\Discharge{!E}{y}$}
\shortDeduce
\DeduceC{$\lforall[x][!B(x)]$}
\DischargeRule{\Elim{\lor}}{x\,y}
\TrinaryInfC{$\lforall[x][!B(x)]$}
\RightLabel{\Elim{\lforall}}
\UnaryInfC{$!B(t)$}
\DisplayProof
\quad$\longrightarrow$\\[8ex]
\quad\AxiomC{$!D \lor !E$}
\AxiomC{$\Discharge{!D}{x}$}
\shortDeduce
\DeduceC{$\lforall[x][!B(x)]$}
\RightLabel{\Elim{\lforall}}
\UnaryInfC{$!B(t)$}
\AxiomC{$\Discharge{!E}{y}$}
\shortDeduce
\DeduceC{$\lforall[x][!B(x)]$}
\RightLabel{\Elim{\lforall}}
\UnaryInfC{$!B(t)$}
\DischargeRule{\Elim{\lor}}{x\,y}
\TrinaryInfC{$!B(t)$}
\DisplayProof
\\
\caption{Beberapa konversi permutasi bagi \Log{N1i}.}
\ollabel{tab:permutation-conversions}
\end{longtable}
}

\begin{prob}
Berikan konversi permutasi bagi \Elim{\lor} yang diikuti \Elim{\lor} atau
\Elim{\lnot}.
\end{prob}

\begin{prob}
Berikan konversi permutasi bagi \Elim{\lexists} yang diikuti
\Elim{\lexists}. Jelaskan mengapa tidak ada syarat variabel eigen yang
dilanggar pada kasus terakhir.
\end{prob}

Kita perlu meyakinkan diri bahwa penerapan konversi permutasi tidak mungkin
melanggar suatu syarat variabel eigen. Pertimbangkan
\[
\AxiomC{}
\RightLabel{$\delta_2$}
\DeduceC{$!D \lor !E$}
\AxiomC{$\Discharge{!D}{x}$}
\RightLabel{$\delta_3$}
\DeduceC{$\lexists[x][!B(x)]$}
\AxiomC{$\Discharge{!E}{y}$}
\RightLabel{$\delta_4$}
\DeduceC{$\lexists[x][!B(x)]$}
\DischargeRule{\Elim{\lor}}{x\,y}
\TrinaryInfC{$\lexists[x][!B(x)]$}
\AxiomC{$\Discharge{!B(c)}{z}$}
\RightLabel{$\delta_5$}
\DeduceC{$!C$}
\DischargeRule{\Elim{\lexists}}{z}
\BinaryInfC{$!C$}
\DisplayProof
\]
Jika kita menggantinya dengan
\[
\AxiomC{}
\RightLabel{$\delta_2$}
\DeduceC{$!D \lor !E$}
\AxiomC{$\Discharge{!D}{x}$}
\RightLabel{$\delta_3$}
\DeduceC{$\lexists[x][!B(x)]$}
\AxiomC{$\Discharge{!B(c)}{z}$}
\RightLabel{$\delta_5$}
\DeduceC{$!C$}
\DischargeRule{\Elim{\lexists}}{z}
\BinaryInfC{$!C$}
\AxiomC{$\Discharge{!E}{y}$}
\RightLabel{$\delta_4$}
\DeduceC{$\lexists[x][!B(x)]$}
\AxiomC{$\Discharge{!B(c)}{z}$}
\RightLabel{$\delta_5$}
\DeduceC{$!C$}
\DischargeRule{\Elim{\lexists}}{z}
\BinaryInfC{$!C$}
\DischargeRule{\Elim{\lor}}{x\,y}
\TrinaryInfC{$!C$}
\DisplayProof
\]
kita mungkin khawatir bahwa variabel eigen~$c$ muncul, misalnya, dalam~$!D$.
Bagaimanapun, asumsi $!D$ dilepaskan di atas \Elim{\lexists} semula, sehingga
dalam !!{proof} semula asumsi itu tidak muncul dalam asumsi yang terbuka di
atas premis utama inferensi \Elim{\lexists}; dengan demikian, syarat
variabel eigen terpenuhi di sana. Namun, dalam !!{proof} baru, $!D$ belum
dilepaskan sampai setelah inferensi-inferensi \Elim{\lexists}, sehingga
syarat variabel eigen akan dilanggar. Akan tetapi, ingatlah bahwa kita
mengandaikan !!{proof}s kita teratur: setiap variabel eigen merupakan
variabel eigen dari tepat satu inferensi dan hanya muncul di atas premis
minor \Elim{\lexists} dalam !!{proof}. Khususnya, $c$ tidak mungkin muncul
dalam $\delta_3$ maupun~$\delta_4$.

Contoh ini menunjukkan hal lain: perhatikan bahwa !!{proof} baru mempunyai
dua salinan~$\delta_5$. Artinya, jika $\delta_5$ memuat suatu potongan
maksimal, kita \emph{belum} mengurangi panjang potongan! Karena itu, kita
harus memastikan bahwa konversi permutasi diterapkan hanya jika kita tahu
tidak ada potongan maksimal dalam $\delta_5$. Kita akan menjaminnya dengan
selalu memilih potongan maksimal yang \emph{paling atas}, yaitu potongan
maksimal sedemikian sehingga sub-!!{proof} (dalam kasus ini,~$\delta_1$)
tidak memuat potongan maksimal yang berakhir di atas inferensi terakhir
dalam sub-!!{proof} tersebut. Hal ini meniadakan kemungkinan adanya
potongan maksimal lain dalam~$\delta_5$ (atau bahkan dalam $\delta_2$,
$\delta_3$, maupun~$\delta_4$).

\end{document}
